% -----------------------------------------------------------------------
% TWIN-LOCATION NOTICE. This file is one of the shared homes for the same
% results; edit every listed home or the drift checker fails.
%   Standalone paper : docs/harbor-research/tex/paper7.tex (S4 mechanism)
%   Chapter          : website-v2/public/whitepaper/
%                       federated-harbor-whitepaper.tex (Chapter 7, The
%                       Federated Harbor; S(fh-sheaf))
%   Index ids        : R6
%   Check            : python3 scripts/harbor-research/check_library_index.py
%                      python3 scripts/harbor-research/check_propagated_corrections.py
%                      python3 scripts/harbor-research/check_citations.py
%   System document  : docs/harbor-research/LIBRARY-SYSTEM.md
% -----------------------------------------------------------------------
\documentclass[11pt,a4paper]{article}
\usepackage[margin=1in]{geometry}
\usepackage{amsmath,amssymb,amsthm}
\usepackage{xcolor}
\usepackage{booktabs}
\usepackage{longtable}
\usepackage{enumitem}
\IfFileExists{titlesec.sty}{%
  \usepackage{titlesec}%
  \titleformat{\section}{\large\bfseries\color{harborblue}}{\thesection}{0.7em}{}%
  \titleformat{\subsection}{\normalsize\bfseries}{\thesubsection}{0.6em}{}%
}{}
\usepackage[colorlinks=true,linkcolor=blue!45!black,urlcolor=blue!45!black]{hyperref}
\usepackage{tikz}
\IfFileExists{pgfplots.sty}{%
  \usepackage{pgfplots}%
  \pgfplotsset{compat=1.18}%
}{}
\usetikzlibrary{arrows.meta,positioning,fit,backgrounds,calc,shapes.geometric,shapes.symbols,decorations.pathmorphing,decorations.pathreplacing}
% --- SYSTEM ARCHITECTURE COLOR PALETTE ---
\definecolor{ptRed}{HTML}{B52227}
\definecolor{ptBlue}{HTML}{274975}
\definecolor{ptDark}{HTML}{1A1A1A}
\definecolor{ptLightGray}{HTML}{F4F4F4}
\definecolor{ptGray}{HTML}{E2E2E2}
\definecolor{ptGold}{HTML}{D4A373}
\definecolor{ptGrid}{HTML}{E5E5E5}
\definecolor{harborblue}{HTML}{274975}
\definecolor{shipred}{HTML}{B52227}
\definecolor{seagreen}{RGB}{31,110,70}
\definecolor{slate}{RGB}{71,85,105}
\definecolor{slatedark}{RGB}{30,41,59}
\definecolor{amber}{HTML}{D4A373}
\definecolor{emerald}{RGB}{16,185,129}
\newtheorem{theorem}{Theorem}
\newtheorem{claim}{Claim}
\theoremstyle{definition}
\newtheorem{change}{Change}
\newtheorem{move}{Move}
\theoremstyle{remark}
\newtheorem*{verdict}{Verdict}
\newtheorem*{adj}{Adjudication}
\setlist{itemsep=1pt,topsep=2pt}

% Shared native-figure styles for the harbor-research corpus (relation-maps,
% regime diagrams, and data plots). Vector TikZ/pgfplots only.
\tikzset{
  canvasgrid/.style={step=0.25cm, draw=ptGrid, line width=0.35pt},
  swissgrid/.style={canvasgrid},
  majorgrid/.style={step=1.00cm, draw=ptGrid!75!black, line width=0.55pt},
  snode/.style={rectangle, draw=ptDark, line width=1pt, fill=white, inner sep=4pt, minimum height=0.65cm, font=\sffamily\bfseries\scriptsize, align=center},
  darknode/.style={snode, fill=ptDark, text=white, font=\sffamily\bfseries\scriptsize\color{white}},
  mask/.style={fill=white, inner sep=1.5pt},
  dArr/.style={-{Stealth[length=2.5mm, width=1.6mm]}, ptDark, line width=1.1pt, shorten >=2pt, shorten <=2pt},
  rArr/.style={-{Stealth[length=2.5mm, width=1.6mm]}, ptRed, line width=1.3pt, shorten >=2pt, shorten <=2pt},
  bArr/.style={-{Stealth[length=2.5mm, width=1.6mm]}, ptBlue, line width=1.1pt, shorten >=2pt, shorten <=2pt},
  relnode/.style={draw,thick,minimum width=2.6cm,minimum height=0.9cm,align=center,font=\small},
  relarrow/.style={-{Stealth[length=2.5mm]},thick},
  regimebox/.style={draw,thick,align=center,font=\small,inner sep=4pt}
}
\IfFileExists{pgfplots.sty}{%
\pgfplotsset{
  harbor curve/.style={
    /pgfplots/width=0.46\textwidth,/pgfplots/height=5.6cm,
    /pgfplots/axis lines=left,
    /pgfplots/label style={font=\scriptsize},
    /pgfplots/tick label style={font=\scriptsize},
    /pgfplots/title style={font=\small,align=center,yshift=2pt},
    /pgfplots/legend style={font=\scriptsize,draw=ptDark!40,line width=0.6pt,fill=white,fill opacity=0.95,text opacity=1,inner sep=3pt},
    thick,
  },
  harbor bar/.style={
    /pgfplots/width=0.46\textwidth,/pgfplots/height=5.6cm,
    /pgfplots/ybar,/pgfplots/bar width=6pt,
    /pgfplots/axis lines=left,
    /pgfplots/label style={font=\scriptsize},
    /pgfplots/tick label style={font=\tiny},
    /pgfplots/xticklabel style={rotate=40,anchor=east,font=\tiny,inner sep=1pt},
    /pgfplots/title style={font=\small,align=center,yshift=2pt},
    /pgfplots/legend style={font=\scriptsize,draw=ptDark!40,line width=0.6pt,fill=white,fill opacity=0.95,text opacity=1,inner sep=3pt},
  },
}
}{}

\usepackage{graphicx}
\graphicspath{{../figures/}}
\usepackage{mdframed}
\newmdenv[backgroundcolor=blue!4,linecolor=harborblue,linewidth=0.8pt,skipabove=6pt,skipbelow=6pt]{thebox}
\newmdenv[backgroundcolor=red!4,linecolor=shipred,linewidth=0.8pt,skipabove=6pt,skipbelow=6pt]{boundary}
\newcommand{\onebreath}[1]{\begin{center}\emph{#1}\end{center}}
\newtheorem{corollary}{Corollary}

\title{\textbf{The Cohomology of Equivocation}\\[4pt]
\large What Cycle Residuals Certify in Federated Witness-Log Gossip\\[3pt]
\normalsize Paper 7 of the Harbor program --- observability, counterexamples, and the consistency radius}
\author{Erich Owens\\ \normalsize The Harbor Research Program \quad \texttt{erich@curiositech.ai}}
\date{August 2026}
\begin{document}
\maketitle

\begin{abstract}
\noindent A cycle residual is a sound test of whether observed edge differences admit one vertex assignment, but its
diagnostic value depends on precisely what evidence the analyst holds. For prefix-coordinate edge differences on a visible
graph $K$, $r=\operatorname{dist}(g_K,\operatorname{im}\delta_K)$ is the norm of the cycle-space projection. A single
nonzero fault on edge $e$ has $r=|s|\sqrt{1-R^{K_c}_{\mathrm{eff}}(e)}$: it is detectable exactly when $e$ lies on a
cycle of the visible coordinate graph $K_c$. This criterion does \emph{not} extend to arbitrary split views by one
witness; a triangle with a leaf supplies an explicit counterexample. For sender $q$, the number of visible-neighbor
components after deleting $q$, minus one, is the exact dimension of invisible split views in each coordinate.
More decisively, if the analyst receives both
authenticated endpoint reports on each relayed link, it can compare a sender's reports directly. Every nonzero cycle
residual then implies an inconsistency that this equal-information checker can already detect. The synthetic harness's
``cohomology-only'' counts compare against a weaker checked-edge baseline and do not establish a detection advantage.
The residual remains useful as a quantitative, topology-aware consistency score when only authenticated edge-difference
receipts are retained, but an exact cycle-sum test has the same noiseless detection set. We give the exact information
boundary, the effective-resistance sensitivity law, and a benchmark contract needed to test whether scoring, localization,
or audit cost yields operational value. No deployed signed-receipt system is evaluated here.
\end{abstract}

\section{What the analyst can know}\label{sec:problem}

A federation sends authenticated witness-log reports along a graph. An auditor wants to know whether one witness gave
incompatible reports to different peers. The evidence question asks which authenticated statements reach the auditor.
The algebraic question asks whether derived edge differences fit one assignment of states to vertices. A cycle residual
answers the second question; the first determines whether that answer adds information.

Suppose witness $q$ sends one log-head claim to $a$ and a different claim to $b$. If both signed claims reach the auditor
with the same source, coordinate, and snapshot identifier, a direct duplicate-claim check detects the conflict. The
check does not require $a$ and $b$ to compare with each other. A detector that discards those claims and retains only
edge differences cannot fairly claim an advantage over the direct check.

A constrained setting still merits study. An auditor may retain authenticated edge-difference receipts while raw
endpoint claims are unavailable. It can then test whether visible cycle sums vanish. The residual gives a magnitude
and a sensitivity map; topology states where the test is blind. The exact evidence contract is part of the result.

\section{Cochains, sections, and cycles}\label{sec:vocab}

A cellular sheaf assigns a vector space to each vertex and edge and a restriction map from each incident vertex into the
edge space. The paper uses coordinate-selection maps. Its real-valued stalks model \emph{features} of log reports,
conditional on an equality-preserving map from authenticated claims; they are not raw cryptographic log heads.

The coboundary $\delta$ maps a vertex assignment $x$ to its oriented edge differences. For one scalar coordinate,
$\delta$ is the signed incidence matrix. The assignment $x$ is a \emph{zero-cochain}; the pattern $\delta x$ is a
\emph{coboundary}. A \emph{global section} of the unshifted sheaf instead satisfies $\delta x=0$. The detector asks
whether the observed edge cochain $g$ lies in $\operatorname{im}\delta$, meaning that some vertex assignment explains
its differences. It does not assert that the explaining assignment is a global section of the unshifted sheaf, or that
the sender claims were mutually consistent.

A signed sum of a coboundary around a graph cycle telescopes to zero. The quotient
$H^1=\operatorname{coker}\delta$ describes edge data modulo coboundaries; in the coordinate-selection graph model,
its Euclidean representative is the projection onto the cycle space. The existence of nonzero $H^1$ says the topology
admits an obstruction. It does not say the observed execution has one: only the class of $g$ can do that.

The completion residual $r$ is the Euclidean distance from $g$ to $\operatorname{im}\delta$ after severed edge rows
are omitted. Effective resistance of an edge is the voltage drop across it under one unit of injected current in the
unit-resistor graph. It governs how strongly a \emph{single-edge} perturbation survives this projection. These objects
can be computed with incidence matrices and least squares; the cohomological language identifies their structural
meaning without adding information to the observed receipts.

\section{The model and the observability contract}\label{sec:model}

\textbf{The sheaf.} Fix a finite gossip graph $G=(V,E)$. Witness $v$ carries a local log state $x_v\in\mathbb{R}^{d_v}$ (its
stalk). Each edge $e=\{u,v\}$ carries the \emph{shared prefix} the two endpoints both report --- a subset $S_e$ of
coordinates, $|S_e|\le L$ --- and the restriction maps $\rho_{v,e}:\mathbb{R}^{d_v}\to\mathbb{R}^{|S_e|}$ are
coordinate-subset selections. The observable on edge $e$ is the disagreement
$g_e=\rho_{u,e}(\text{$u$'s report})-\rho_{v,e}(\text{$v$'s report})$. A signed head requires an explicit
equality-preserving feature map before this real-valued subtraction: equal heads map equally, and a head inequality
claimed detectable in coordinate $c$ must map to an inequality in $c$. No such signed-receipt implementation is evaluated
here. Receipts must bind source, peer, coordinate, snapshot, and provenance before statements can be compared. Honesty means all sender reports are consistent with one vertex assignment
$x$: then $g=\delta x$ is a coboundary of the sheaf coboundary $\delta$, and every cycle-sum of $g$ vanishes. Equivocation is
an offset pattern $\varepsilon$ added to $q$'s per-neighbor reports; the question for the derived edge data is what survives
in $H^1$ --- the obstruction space $\operatorname{coker}(\delta)$, which for
prefix restrictions decomposes per shared coordinate $c$ into the cycle spaces of the \emph{coordinate subgraphs}
\[
G_c \;=\; \bigl(V,\ \{e\in E:\ c\in S_e\}\bigr)
\]
[verified framework, \cite{curry,hansen-ghrist}]. $G_c$ is what the \emph{federation} has: every link that carries
coordinate $c$. What the \emph{edge-only analyst} has is smaller when links are severed; see $K_c$ below.

\textbf{Three-tier visibility --- the observability contract.} Not every edge is equally visible to the analyst, and the
paper's claims are stated \emph{per tier} (Figure~\ref{fig:vis}):
\begin{itemize}
\item \textbf{compared} --- the analyst holds the direct cross-check of the edge: pairwise detection, no cohomology needed;
\item \textbf{relayed} --- both authenticated endpoint claims reach the analyst by gossip, but the endpoints did not
cross-check. Their arrival permits both a direct duplicate-claim check by sender and a derived edge difference;
\item \textbf{severed} --- no report about the edge crosses to the analyst at all: the corresponding data block is a free
variable of the edge-data completion. Other authenticated claims from that sender may still be available elsewhere.
\end{itemize}

\textbf{The visible coordinate subgraph $K_c$.} Write $K$ for the compared-plus-relayed edge set --- every link the analyst
holds a report about --- and
\[
K_c \;=\; \bigl(V,\ \{e\in K:\ c\in S_e\}\bigr)\ \subseteq\ G_c
\]
for its coordinate-$c$ part: the coordinate subgraph with the severed links deleted. Every \emph{one-edge residual} criterion
is a statement about cycles of $K_c$, never of $G_c$, and the reason is one line of algebra made explicit in the detector
below --- a severed block is a free variable, so minimising over it simply deletes its row, and a loop that runs through a
severed link is not a loop the analyst can close. The federation's topology sets $G_c$; the observability contract sets
$K_c$; only the second enters this edge-data residual.

\textbf{The detector.} The \emph{completion residual} is the least-squares defect of the best global explanation of
everything visible:
\[
r \;=\; \min_{x,\;g_{\mathrm{sev}}}\ \bigl\|\,g_K-(\delta x)\big|_K\,\bigr\|_2,
\]
the minimum over vertex assignments $x$ and free severed blocks $g_{\mathrm{sev}}$, where $K$ is the compared-plus-relayed
edge set. The severed blocks are wholly unconstrained, so minimising over them is exactly \emph{dropping their rows}:
$g_{\mathrm{sev}}$ never enters the objective, and $r=\operatorname{dist}\bigl(g_K,\ \operatorname{im}\delta_K\bigr)$ for
$\delta_K$ the coboundary restricted to visible edges. Honest executions give $r=0$ exactly, for every visibility pattern;
$r>0$ certifies that \emph{no} global assignment explains the visible edge differences. Its size depends on the feature
map, coordinate scaling, and metric; it is not a cryptographic measure of log divergence. Equivalently $r=\|\Pi_K\,g_K\|$, the projection of
the visible disagreement data onto the per-coordinate cycle spaces \emph{of the visible subgraphs}:
\[
r^2=\sum_c\bigl\|\mathrm{Proj}_{\mathrm{cycle}(K_c)}\,g^c\bigr\|^2,
\qquad g^c\in\mathbb{R}^{K_c}\ \text{the coordinate-$c$ part of the visible data}
\]
[internal, \texttt{sheaf\_harness\_v2.py}]. The derivation is three lines and is the whole bridge from the minimization to
Theorem 1, so it is worth stating: prefix restrictions are coordinate selections, so $\delta_K$ block-diagonalises per
coordinate into the signed incidence operator $B_c$ of $K_c$ --- the severed rows having already dropped out; edge space
splits orthogonally as $\mathbb{R}^{K_c}=\operatorname{im}(B_c)\oplus\ker(B_c^{\!\top})$, cut space plus cycle space; and the
residual is what is left after removing everything an honest world could have produced, i.e.\ the component in
$\mathrm{cycle}(K_c)$. Writing $G_c$ here instead would not merely be loose, it would fabricate detections: on $C_6$ with a
lie of $3$ on a relayed edge and one other cycle edge severed, the residual is $1.1\times10^{-14}$ while the $G_c$ form
reports $1.2247$ --- and reports a different value again for a different honest world, so it is not even a well-defined
statistic [internal, \texttt{sheaf\_consistency\_radius.py}].

\input{fig-paper7-visibility.tex}

\section{The mechanism: visible cycles and their exact scope}\label{sec:mechanism}

\onebreath{A one-edge fault is visible to the residual exactly when its edge lies on a visible coordinate cycle. A general split view by one sender obeys a different criterion, and authenticated relayed claims admit direct comparison.}

\begin{thebox}
\textbf{Theorem 1 (one-edge criterion).} Fix a coordinate $c$, a visible edge $e$ of $K_c$, and a nonzero edge perturbation $s$. With unit edge weights and the Euclidean metric,
\[
r=|s|\sqrt{1-R_{\mathrm{eff}}^{K_c}(e)}.
\]
Consequently $r>0$ exactly when $e$ lies on a cycle of $K_c$. A bridge has effective resistance $1$ and zero residual. A severed edge contributes no observed row. This is a theorem about a \emph{one-edge} perturbation, not every split view by a single sender.

\textbf{General one-sender criterion.} Let $A_qo$ place sender $q$'s per-neighbor coordinate offsets on visible incident edges, with orientation signs. Then $r=\|\Pi_K A_qo\|$, so detection is exactly $\Pi_K A_qo\ne0$. For each coordinate, $\ker(\Pi_KA_q)$ comprises offset patterns constant on edges from $q$ into each connected component of $K_c-q$. Constants for distinct components may differ. Indeed, $\Pi_KA_qo=0$ exactly when $A_qo=B_cx$ for a potential $x$. Every nonincident edge has zero offset, forcing $x$ to be constant on each component of $K_c-q$; those constants also construct the converse potential.

\textbf{Blind-spot dimension.} Let $d_q^c$ be the number of visible $c$-edges at $q$, and let $h_q^c$ be the number of connected components of $K_c-q$ touched by those edges. If $d_q^c>0$, then
\[
\dim\ker(\Pi_KA_q)^c=h_q^c,\qquad
b_q^c:=\dim\bigl(\ker(\Pi_KA_q)^c/\langle\mathbf1\rangle\bigr)=h_q^c-1.
\]
Set $b_q^c=0$ when $d_q^c=0$. The quotient removes a uniform sender-state change, leaving precisely the independent \emph{invisible split-view} directions. Thus every nonuniform one-sender offset in coordinate $c$ is detected exactly when $b_q^c=0$. The total blind-spot dimension for independent coordinates is $B_q=\sum_c b_q^c$. Under an edge-only receipt design that may add $c$-bearing edges among the touched components of $K_c-q$, at least $b_q^c$ such edges are needed to remove this blindness, and $b_q^c$ suffice if arbitrary cross-component links are admissible. Each added edge can merge at most two components; a spanning tree on the components attains the bound. This counts observability, not authenticated-claim consistency or Byzantine fault tolerance.
\end{thebox}

\textbf{Counterexample to a universal split-view criterion.} Let $q,a,b$ form a triangle and attach leaf $l$ to $q$, all edges visible in one coordinate. Starting from zero states, $q$ sends offset $1$ to $a,b$ and $0$ to $l$. The alternative potential $x_q=1$, $x_a=x_b=0$, $x_l=1$ produces exactly these edge differences: the two triangle edges at $q$ differ by $1$, while $a-b$ and $q-l$ differ by $0$. Hence $r=0$ despite a split view and a visible cycle through $q$. The components of $K_c-q$ are $\{a,b\}$ and $\{l\}$, so this is precisely a nonuniform kernel offset. A duplicate-claim checker with both authenticated claims from $q$ detects it directly if they bind the same coordinate and snapshot.

\textbf{The cochain by hand.} On a visible four-relay ring in one scalar coordinate, take the oriented edge differences $g=(0,0,0,3)$. A coboundary sums to zero around the ring; $g$ sums to $3$. The nearest potential $x=(0,0.75,1.5,2.25)$ gives $\delta x=(-0.75,-0.75,-0.75,2.25)$, leaving magnitude $0.75$ on each edge. Thus
\[
r=\sqrt{4(0.75)^2}=1.5=3/\sqrt4=3\sqrt{1-R_{\mathrm{eff}}^{C_4}(e)},\quad R_{\mathrm{eff}}^{C_4}(e)=3/4.
\]
On a path, or on the ring with any other edge severed, the same one-edge fault is a coboundary and $r=0$. On a fully visible $C_6$, a fault of size $3$ gives $3/\sqrt6=1.2247$; severing another edge makes the residual zero. These exact scalar examples preserve the valid effective-resistance mechanism.

\input{fig-r6-regime.tex}

\textbf{Equal-information baseline.} In the stated relayed tier the analyst holds authenticated endpoint reports. It can group a sender's claims by coordinate and snapshot and test equality without asking the endpoints to compare. If every sender made one consistent claim, all derived edge differences would be a coboundary; therefore any nonzero cycle residual implies an inconsistency that this direct checker can also detect. An edge-only receipt pipeline may lose the raw claims; there the residual still gives a metric score and residual pattern. With the same edge-only data, exact oriented cycle sums have the identical noiseless zero/nonzero detection set. The residual does not create additional detection information.

\begin{thebox}
\textbf{Proposition 2 (equal-information dominance).} Assume each visible edge record contains both authenticated endpoint
claims for the same coordinate and snapshot, and the edge cochain is their difference. If a sender's claims agree across
all its visible incident edges for each coordinate, then $g_K=\delta_Kx$ for a vertex assignment $x$ and $r=0$.
Consequently a nonzero residual implies a conflict that direct duplicate-claim checking on the \emph{same records}
detects. The converse fails, including in the triangle-with-leaf example.

\emph{Proof.} Define $x_v^c$ as sender $v$'s common claim for coordinate $c$; choose any value when $v$ has no visible
incident edge carrying $c$. Every visible edge difference is then $x_u^c-x_v^c$, so $g^c=B_cx^c$ and its cycle
projection vanishes. Contraposition proves the detection implication.
\end{thebox}

\begin{boundary}
A zero residual means the visible edge differences admit a potential, not that the actual reports came from one head. A nonzero residual proves inconsistency of the encoded, snapshot-bound data, not who caused it. Both interpretations require authenticated provenance and a feature map that preserves the head equality being tested.
\end{boundary}

\section{The harness: what its trials establish}\label{sec:harness}

The synthetic harness uses 200 trials per arm at seed 20260816, real five-dimensional stalks, and shared coordinate
subsets. The checked-edge baseline tests whether a directly checked edge exposed the fault. A trial labeled
cohomology-only means the residual was positive while this restricted baseline was silent. The harness also
reconstructs each sender's visible coordinate claims and runs Proposition~2's equal-information duplicate-claim check.

\begin{thebox}
\textbf{Observed behavior.} In the two-path split-view arm the residual was positive and the checked-edge baseline silent
in $200/200$ trials. In the expander arm the corresponding count was $113/200$, with $87$ residual-dark trials. Of those
$87$, $79$ had no relayed incident fault edge and $8$ had a relayed fault edge on no visible coordinate cycle. The
severed-cycle arm gave $0/200$; a cut-edge arm had maximum numerical residual $1.5\times10^{-13}$ over 400 trials. The
no-equivocator control gave $0/200$. Across all nine arms and 1,800 trials, the equal-information checker found
\emph{zero} residual-only cases and \emph{209} direct-only cases. These counts support Proposition~2 within the
synthetic feature and snapshot model; they are not a deployed signed-log evaluation.
\end{thebox}

The counts depend on the sampled topology, offsets, metric, and threshold. The harness checks the blind-spot dimension against independent sender-map nullity calculations in all 224 four-vertex graph/sender cases with at least one incident edge. The equal-information comparison has been
run on these synthetic records; exact cycle sums and noisy scoring have not been compared. An operational benchmark
must predefine authenticated receipt and snapshot binding, the equality-preserving feature map, noise, false-alarm
threshold, runtime cost, and localization target, then run all three evaluators on matched evidence.

\section{The consistency radius: certificate, sensitivity, and cost}\label{sec:radius}

\onebreath{The residual is the least Euclidean unrestricted edge perturbation needed to reconcile visible data. A perturbation constrained to one sender can require more, or be infeasible. Effective resistance supplies exact one-edge sensitivity.}

\begin{thebox}
\textbf{CR-1 (soundness and constrained lower bound).} Let $\Pi_K$ project onto $\ker(\delta_K^{\!\top})$ in the chosen Euclidean metric. Then
\[
r=\min_x\|g_K-\delta_Kx\|=\|\Pi_Kg_K\|.
\]
Thus $r>0$ exactly when no potential explains the visible differences, and honest data give $r=0$. The minimum norm of an \emph{unrestricted edge-data perturbation} making $g_K$ a coboundary is exactly $r$, achieved by $-\Pi_Kg_K$. If the perturbation must originate at a fixed witness $q$, the feasible minimum is
\[
m_q(g)=\inf\{\|A_qo\|:\Pi_KA_qo=\Pi_Kg_K\}\ge r,
\]
with $m_q(g)=+\infty$ when the constraint is infeasible. The minimum norm of the offset vector $o$ is a separate quantity. With $o_\perp$ orthogonal to $\ker(\Pi_KA_q)$,
\[
\sigma_{\min}^{+}(\Pi_KA_q)\|o_\perp\|
\le \|\Pi_KA_qo\|\le\|\Pi_KA_q\|\|o_\perp\|.
\]
Kernel offsets can be arbitrarily large while $r$ remains zero.

\textbf{CR-2 (one-edge sensitivity).} For a visible unit-weight edge $e$ in coordinate $c$,
\[
\|\Pi_K(s\mathbf e_e)\|=|s|\sqrt{1-R_{\mathrm{eff}}^{K_c}(e)}.
\]
Its projected support lies on visible cycles containing $e$. For general sender offsets, support lies in the union of cycle-containing blocks reached by incident offsets; cancellation can remove some or all signal, and the pattern cannot attribute a sender. Additive edge noise $\eta$ changes the residual by at most $\|\Pi_K\eta\|$. A threshold requires a justified noise bound or calibration.

\textbf{CR-3 (coordinate decomposition).} Let $C=\bigcup_{e\in K}S_e$ and $T=|C|$. Coordinate selection splits the projection into at most $T$ scalar graph-Laplacian systems, one for each nonempty $K_c$. The bound is at most $L$ only if all edge coordinate sets draw from one globally common set of at most $L$ coordinates. Cost depends on $\sum_c|K_c|$, conditioning, tolerance, and solver. Near-linear sparse-solver theory~\cite{spielman-teng} is relevant but does not prove a per-round runtime for this federation.
\end{thebox}

\textbf{Numbers by hand.} On $C_n$, the edge is in parallel with a path of $n-1$ unit resistors, so $R_{\mathrm{eff}}(e)=(n-1)/n$ and a one-edge size-$s$ perturbation gives $r=|s|/\sqrt n$. On $C_6$, the projected norm per unit one-edge fault is $1/\sqrt6\approx0.4082$ and $s=3$ yields $1.2247$; on $C_{12}$, $s=0.5$ yields about $0.1443$. This $0.4082$ is a one-edge sensitivity, not a universal minimum singular value for arbitrary sender offsets. These are magnitudes in a selected real-coordinate metric, not cryptographic measures of signed-head divergence. The harness's achievability construction realizes an unrestricted edge perturbation; it does not show that a fixed sender can attain the same minimum.

\input{fig-paper7-radius.tex}

Effective resistance maps where a one-edge fault is hard to see in edge-only data. Adding independent visible
detours can lower the effective resistance of an edge and increase its one-edge sensitivity. Direct comparison remains
the equal-information detection baseline, while scoring, localization, and audit cost need a common benchmark.

\section{Related work and contribution}\label{sec:related}

Cellular sheaves on graphs and their coboundary calculus follow Curry~\cite{curry}; spectral sheaf methods follow
Hansen--Ghrist~\cite{hansen-ghrist}. HodgeRank already decomposes edge observations into a best-fit gradient and a
cyclic residual~\cite{hodgerank}. Effective resistance as an incidence-row leverage score is also established
~\cite{drineas-mahoney}. The language of local consistency without a global assignment is related to
Abramsky--Brandenburger contextuality~\cite{abramsky-brandenburger}. Robinson~\cite{robinson} supplies the
consistency-radius perspective. The projection, cycle-space identity, and effective-resistance law here are familiar
graph linear algebra rather than a new cohomology invariant.

PeerReview~\cite{peerreview} and BFT forensics~\cite{sheng} motivate preserving authenticated statements that can identify a conflicting sender. When both endpoint claims are available to the analyst, a direct duplicate-claim check has that evidentiary advantage. The present real-coordinate model instead characterizes the consistency score of derived edge differences. An exact cycle-sum checker sees the same noiseless inconsistency as its least-squares residual; the residual adds a metric magnitude and a topology-dependent distribution of error. The contribution is an explicit information contract, the one-edge sensitivity formula under visible-edge deletion, the general single-sender kernel characterization, and a counterexample that prevents extending the one-edge criterion to arbitrary split views.

\section{Limits of the certificate}\label{sec:boundary}

\begin{boundary}
\begin{itemize}
\item A zero residual certifies only the existence of a potential for visible encoded edge differences. The triangle-with-leaf example is a one-sender split view with zero residual. Separate sender claims may still expose it. A coalition can also arrange canceling edge offsets.
\item A nonzero residual does not name a culprit. Different signed-report histories can yield the same edge-difference vector; retaining authenticated claims is necessary for direct attribution.
\item A severed edge contributes no evidence. Deleting an edge can also break the only visible cycle supporting a one-edge residual. This is a statement about the specified edge-data pipeline, not a claim that all evidence about that witness is absent.
\item The numerical radius depends on the equality-preserving feature map and metric. With measurement noise $\eta$, the honest residual may be as large as $\|\Pi_K\eta\|$; operational thresholds need a noise model or calibration. A zero real-coordinate obstruction does not establish consistency in every richer coefficient system~\cite{caru}.
\end{itemize}
\end{boundary}

\section*{Reproducibility}

The synthetic scripts \texttt{sheaf\_mechanism\_proof.py}, \texttt{sheaf\_harness\_v2.py}, and
\texttt{sheaf\_consistency\_radius.py} use seed 20260816 for their sampled arms. They reproduce the scalar cycle
examples, visibility-mask checks, and effective-resistance calculations. The harness evaluates both the restricted
checked-edge baseline and an equal-information duplicate-claim checker; an assertion rejects any residual-only case
against the latter. Its unrestricted completion construction does not solve the fixed-sender minimum. Its timing and
localization reports depend on the sampled graphs, metric, and machine; no deployed signed-receipt system is tested.
The triangle-with-leaf counterexample and the componentwise kernel proof above are independent hand checks.

\begin{thebibliography}{9}
\bibitem{peerreview} A.~Haeberlen, P.~Kouznetsov, and P.~Druschel. PeerReview: practical accountability for distributed
systems. In \emph{Proc.\ 21st ACM Symposium on Operating Systems Principles (SOSP)}, pp.~175--188, 2007.
\bibitem{abramsky-brandenburger} S.~Abramsky and A.~Brandenburger. The sheaf-theoretic structure of non-locality and
contextuality. \emph{New Journal of Physics}, 13:113036, 2011.
\bibitem{curry} J.~M. Curry. \emph{Sheaves, Cosheaves and Applications}. PhD thesis, University of Pennsylvania, 2014.
\bibitem{hansen-ghrist} J.~Hansen and R.~Ghrist. Toward a spectral theory of cellular sheaves. \emph{Journal of Applied
and Computational Topology}, 3:315--358, 2019.
\bibitem{hodgerank} X.~Jiang, L.-H.~Lim, Y.~Yao, and Y.~Ye. Statistical ranking and combinatorial Hodge theory.
\emph{Mathematical Programming}, 127:203--244, 2011.
\bibitem{drineas-mahoney} P.~Drineas and M.~W.~Mahoney. Effective resistances, statistical leverage, and applications
to linear equation solving. arXiv:1005.3097, 2010.
\bibitem{robinson} M.~Robinson. Sheaves are the canonical data structure for sensor integration. \emph{Information
Fusion}, 36:208--224, 2017.
\bibitem{caru} G.~Car\`u. On the cohomology of contextuality. In \emph{Proc.\ 13th International Conference on Quantum
Physics and Logic (QPL 2016)}, EPTCS 236, pp.~21--39, 2017.
\bibitem{sheng} P.~Sheng, G.~Wang, K.~Nayak, S.~Kannan, and P.~Viswanath. BFT protocol forensics. In \emph{Proc.\ ACM
Conference on Computer and Communications Security (CCS)}, pp.~1722--1743, 2021.
\bibitem{spielman-teng} D.~A. Spielman and S.-H. Teng. Nearly-linear time algorithms for graph partitioning, graph
sparsification, and solving linear systems. In \emph{Proc.\ 36th ACM Symposium on Theory of Computing (STOC)},
pp.~81--90, 2004.
\end{thebibliography}

\end{document}
